% --- Archon physics formalization source begin ---
% archon:physics
% archon:covers PhyXMiniProblems/problem_phyx_mini_0659.lean
% archon:source-report reports/phyx_mini/problem_phyx_mini_0659.source.json

\chapter{Physics problem phyx\_mini\_0659}
\label{ch:PhyXMiniProblems_problem_phyx_mini_0659}

\paragraph{Problem source.}
\#\# Physical scenario

Figure shows a few energy levels of the mercury atom.

\#\# Question

What minimum speed must an electron have to excite the \$492\$-\$nm\$-wavelength blue emission line in the \$Hg\$ spectrum?

\#\# Answer choices

- A: A: \$1.40 \textbackslash{}times 10\textasciicircum{}6\$ \$m/s\$

- B: B: \$1.60 \textbackslash{}times 10\textasciicircum{}6\$ \$m/s\$

- C: C: \$1.80 \textbackslash{}times 10\textasciicircum{}6\$ \$m/s\$

- D: D: \$2.0 \textbackslash{}times 10\textasciicircum{}6\$ \$m/s\$

\#\# Figure caption (auxiliary; use the image as primary evidence)

The image is an energy level diagram with the y-axis labeled "Energy (eV)" ranging from 0 to 10. Several horizontal lines represent energy levels with their corresponding numerical values in electron volts (eV) and electron configurations. 

1. **At 0.00 eV**: The level is labeled \textbackslash{}(6s\textasciicircum{}2\textbackslash{}).
2. **At 6.70 eV**: The level is labeled \textbackslash{}(6s6p\textbackslash{}).
3. **At 7.92 eV**: The level is labeled \textbackslash{}(6s7s\textbackslash{}).
4. **At 8.84 eV**: Two levels are labeled \textbackslash{}(6s7p\textbackslash{}) and \textbackslash{}(6s6d\textbackslash{}).
5. **At 9.22 eV**: The level is labeled \textbackslash{}(6s8s\textbackslash{}).
6. **At 9.53 eV**: The level is labeled \textbackslash{}(6s8p\textbackslash{}).

At the bottom right, there's text stating, "Energies for each level are in eV" in blue. Horizontal lines are shown under the energy levels to separate them visually.

\#\# Dataset metadata

Quantum Phenomena; Physical Model Grounding Reasoning, Multi-Formula Reasoning

\paragraph{Recorded answer/context.}
C: C: \$1.80 \textbackslash{}times 10\textasciicircum{}6\$ \$m/s\$

\paragraph{Figure/image path.}
/root/proposal\_for\_physic/science-mango/phyx\_mini\_run/phyx\_data/test\_image/659.png

\paragraph{Formalization target.}
create a compiling Lean file with sorry bodies at `PhyXMiniProblems/problem\_phyx\_mini\_0659.lean`. The Lean declarations must preserve the physical quantities, dimensions or dimensional roles, figure labels, governing-law hypotheses, and final relation expressed by this problem.
Use Mathlib/Physlib names found through LeanExplore where available. If a physics API is missing, introduce faithful local abstractions rather than scalar placeholder aliases.

\begin{theorem}[Physics formalization target]
\label{thm:physics:phyx_mini_0659:target}
\lean{PhyXMiniProblems.ProblemPhyXMini0659.minimumElectronSpeed_for_492nmMercuryEmission}
\uses{def:physics:phyx-mini-0659:phyxminiproblems-problemphyxmini0659-mercuryexcitationsetup, def:physics:phyx-mini-0659:phyxminiproblems-problemphyxmini0659-matchesmercuryelectronexcitationscenario, def:physics:phyx-mini-0659:phyxminiproblems-problemphyxmini0659-matchessuppliedmercuryenergyleveldiagram, def:physics:phyx-mini-0659:phyxminiproblems-problemphyxmini0659-usesstandardelectronandphotonreferencedata, def:physics:phyx-mini-0659:phyxminiproblems-problemphyxmini0659-hasphysicalmercuryexcitationparameters, def:physics:phyx-mini-0659:phyxminiproblems-problemphyxmini0659-satisfiesmercuryemissionandthresholdlaws, def:physics:phyx-mini-0659:phyxminiproblems-problemphyxmini0659-agreeswithdisplayedminimumspeed, def:physics:phyx-mini-0659:phyxminiproblems-problemphyxmini0659-isuniqueclosestminimumspeedchoice}
The assigned autoformalize agent should translate this physics problem into Lean declarations in the covered file, with theorem and lemma proof bodies written as `by sorry`.
\end{theorem}
\begin{proof}
This is an autoformalization task, not a proof task. Produce statements that can later be proved by the physics prover without weakening the physical model.
\end{proof}

% --- Archon declaration topology begin ---
\section{Declaration topology}
\begin{definition}[Length Quantity]
  \label{def:physics:phyx-mini-0659:phyxminiproblems-problemphyxmini0659-lengthquantity}
  \lean{PhyXMiniProblems.ProblemPhyXMini0659.LengthQuantity}
  A nonnegative, unit-independent physical length
\end{definition}
\begin{proof}
  This is the declared model definition of Length Quantity.
\end{proof}
\begin{definition}[Mass Quantity]
  \label{def:physics:phyx-mini-0659:phyxminiproblems-problemphyxmini0659-massquantity}
  \lean{PhyXMiniProblems.ProblemPhyXMini0659.MassQuantity}
  A nonnegative, unit-independent physical mass
\end{definition}
\begin{proof}
  This is the declared model definition of Mass Quantity.
\end{proof}
\begin{definition}[action Dimension]
  \label{def:physics:phyx-mini-0659:phyxminiproblems-problemphyxmini0659-actiondimension}
  \lean{PhyXMiniProblems.ProblemPhyXMini0659.actionDimension}
  The physical dimension of action, mass * length\textasciicircum{}2 / time
\end{definition}
\begin{proof}
  This is the declared model definition of action Dimension.
\end{proof}
\begin{definition}[Action Quantity]
  \label{def:physics:phyx-mini-0659:phyxminiproblems-problemphyxmini0659-actionquantity}
  \lean{PhyXMiniProblems.ProblemPhyXMini0659.ActionQuantity}
  \uses{def:physics:phyx-mini-0659:phyxminiproblems-problemphyxmini0659-actiondimension}
  A nonnegative, unit-independent physical action
\end{definition}
\begin{proof}
  This is the declared model definition of Action Quantity.
\end{proof}
\begin{definition}[Light Speed Quantity]
  \label{def:physics:phyx-mini-0659:phyxminiproblems-problemphyxmini0659-lightspeedquantity}
  \lean{PhyXMiniProblems.ProblemPhyXMini0659.LightSpeedQuantity}
  A signed physical speed, used for the vacuum speed of light
\end{definition}
\begin{proof}
  This is the declared model definition of Light Speed Quantity.
\end{proof}
\begin{definition}[length Readout]
  \label{def:physics:phyx-mini-0659:phyxminiproblems-problemphyxmini0659-lengthreadout}
  \lean{PhyXMiniProblems.ProblemPhyXMini0659.lengthReadout}
  \uses{def:physics:phyx-mini-0659:phyxminiproblems-problemphyxmini0659-lengthquantity}
  Read a physical length in a selected unit
\end{definition}
\begin{proof}
  This is the declared model definition of length Readout.
\end{proof}
\begin{definition}[length In Meters]
  \label{def:physics:phyx-mini-0659:phyxminiproblems-problemphyxmini0659-lengthinmeters}
  \lean{PhyXMiniProblems.ProblemPhyXMini0659.lengthInMeters}
  \uses{def:physics:phyx-mini-0659:phyxminiproblems-problemphyxmini0659-lengthquantity, def:physics:phyx-mini-0659:phyxminiproblems-problemphyxmini0659-lengthreadout}
  Metre readout of a physical length
\end{definition}
\begin{proof}
  This is the declared model definition of length In Meters.
\end{proof}
\begin{definition}[length In Nanometers]
  \label{def:physics:phyx-mini-0659:phyxminiproblems-problemphyxmini0659-lengthinnanometers}
  \lean{PhyXMiniProblems.ProblemPhyXMini0659.lengthInNanometers}
  \uses{def:physics:phyx-mini-0659:phyxminiproblems-problemphyxmini0659-lengthquantity, def:physics:phyx-mini-0659:phyxminiproblems-problemphyxmini0659-lengthreadout}
  Nanometre readout of a physical length
\end{definition}
\begin{proof}
  This is the declared model definition of length In Nanometers.
\end{proof}
\begin{definition}[mass In Kilograms]
  \label{def:physics:phyx-mini-0659:phyxminiproblems-problemphyxmini0659-massinkilograms}
  \lean{PhyXMiniProblems.ProblemPhyXMini0659.massInKilograms}
  \uses{def:physics:phyx-mini-0659:phyxminiproblems-problemphyxmini0659-massquantity}
  Kilogram readout of a physical mass
\end{definition}
\begin{proof}
  This is the declared model definition of mass In Kilograms.
\end{proof}
\begin{definition}[energy In Joules]
  \label{def:physics:phyx-mini-0659:phyxminiproblems-problemphyxmini0659-energyinjoules}
  \lean{PhyXMiniProblems.ProblemPhyXMini0659.energyInJoules}
  Coherent-SI readout of a physical energy, in joules
\end{definition}
\begin{proof}
  This is the declared model definition of energy In Joules.
\end{proof}
\begin{definition}[energy In Electron Volts]
  \label{def:physics:phyx-mini-0659:phyxminiproblems-problemphyxmini0659-energyinelectronvolts}
  \lean{PhyXMiniProblems.ProblemPhyXMini0659.energyInElectronVolts}
  \uses{def:physics:phyx-mini-0659:phyxminiproblems-problemphyxmini0659-energyinjoules}
  Electron-volt readout grounded by DimEnergy.electronVolt
\end{definition}
\begin{proof}
  This is the declared model definition of energy In Electron Volts.
\end{proof}
\begin{definition}[action In Joule Seconds]
  \label{def:physics:phyx-mini-0659:phyxminiproblems-problemphyxmini0659-actioninjouleseconds}
  \lean{PhyXMiniProblems.ProblemPhyXMini0659.actionInJouleSeconds}
  \uses{def:physics:phyx-mini-0659:phyxminiproblems-problemphyxmini0659-actionquantity}
  Coherent-SI readout of an action, in joule-seconds
\end{definition}
\begin{proof}
  This is the declared model definition of action In Joule Seconds.
\end{proof}
\begin{definition}[speed In Meters Per Second]
  \label{def:physics:phyx-mini-0659:phyxminiproblems-problemphyxmini0659-speedinmeterspersecond}
  \lean{PhyXMiniProblems.ProblemPhyXMini0659.speedInMetersPerSecond}
  Metres-per-second readout of a nonnegative physical speed
\end{definition}
\begin{proof}
  This is the declared model definition of speed In Meters Per Second.
\end{proof}
\begin{definition}[light Speed In Meters Per Second]
  \label{def:physics:phyx-mini-0659:phyxminiproblems-problemphyxmini0659-lightspeedinmeterspersecond}
  \lean{PhyXMiniProblems.ProblemPhyXMini0659.lightSpeedInMetersPerSecond}
  \uses{def:physics:phyx-mini-0659:phyxminiproblems-problemphyxmini0659-lightspeedquantity}
  Metres-per-second readout of a signed physical speed such as c
\end{definition}
\begin{proof}
  This is the declared model definition of light Speed In Meters Per Second.
\end{proof}
\begin{definition}[Atomic Species]
  \label{def:physics:phyx-mini-0659:phyxminiproblems-problemphyxmini0659-atomicspecies}
  \lean{PhyXMiniProblems.ProblemPhyXMini0659.AtomicSpecies}
  Atomic species distinguished by the problem statement
\end{definition}
\begin{proof}
  This is the declared model definition of Atomic Species.
\end{proof}
\begin{definition}[Projectile Species]
  \label{def:physics:phyx-mini-0659:phyxminiproblems-problemphyxmini0659-projectilespecies}
  \lean{PhyXMiniProblems.ProblemPhyXMini0659.ProjectileSpecies}
  Species of the projectile which excites the atom
\end{definition}
\begin{proof}
  This is the declared model definition of Projectile Species.
\end{proof}
\begin{definition}[Mercury Energy Level]
  \label{def:physics:phyx-mini-0659:phyxminiproblems-problemphyxmini0659-mercuryenergylevel}
  \lean{PhyXMiniProblems.ProblemPhyXMini0659.MercuryEnergyLevel}
  The seven configuration-labelled horizontal levels visible in image 659.png. The separate 6s7p and 6s6d constructors preserve the twofold 8.84 eV readout rather than identifying the physical states
\end{definition}
\begin{proof}
  This is the declared model definition of Mercury Energy Level.
\end{proof}
\begin{definition}[Spectral Color]
  \label{def:physics:phyx-mini-0659:phyxminiproblems-problemphyxmini0659-spectralcolor}
  \lean{PhyXMiniProblems.ProblemPhyXMini0659.SpectralColor}
  Color name assigned to the spectral line in the question
\end{definition}
\begin{proof}
  This is the declared model definition of Spectral Color.
\end{proof}
\begin{definition}[Figure Axis Quantity]
  \label{def:physics:phyx-mini-0659:phyxminiproblems-problemphyxmini0659-figureaxisquantity}
  \lean{PhyXMiniProblems.ProblemPhyXMini0659.FigureAxisQuantity}
  Physical quantity represented by the sole plotted axis
\end{definition}
\begin{proof}
  This is the declared model definition of Figure Axis Quantity.
\end{proof}
\begin{definition}[Figure Energy Unit]
  \label{def:physics:phyx-mini-0659:phyxminiproblems-problemphyxmini0659-figureenergyunit}
  \lean{PhyXMiniProblems.ProblemPhyXMini0659.FigureEnergyUnit}
  Unit printed beside the energy-axis label
\end{definition}
\begin{proof}
  This is the declared model definition of Figure Energy Unit.
\end{proof}
\begin{definition}[Mercury Energy Level Diagram]
  \label{def:physics:phyx-mini-0659:phyxminiproblems-problemphyxmini0659-mercuryenergyleveldiagram}
  \lean{PhyXMiniProblems.ProblemPhyXMini0659.MercuryEnergyLevelDiagram}
  \uses{def:physics:phyx-mini-0659:phyxminiproblems-problemphyxmini0659-mercuryenergylevel, def:physics:phyx-mini-0659:phyxminiproblems-problemphyxmini0659-figureaxisquantity, def:physics:phyx-mini-0659:phyxminiproblems-problemphyxmini0659-figureenergyunit}
  Literal qualitative and numerical content visible in the supplied image
\end{definition}
\begin{proof}
  This is the declared model definition of Mercury Energy Level Diagram.
\end{proof}
\begin{definition}[Mercury Emission Line]
  \label{def:physics:phyx-mini-0659:phyxminiproblems-problemphyxmini0659-mercuryemissionline}
  \lean{PhyXMiniProblems.ProblemPhyXMini0659.MercuryEmissionLine}
  \uses{def:physics:phyx-mini-0659:phyxminiproblems-problemphyxmini0659-lengthquantity, def:physics:phyx-mini-0659:phyxminiproblems-problemphyxmini0659-mercuryenergylevel, def:physics:phyx-mini-0659:phyxminiproblems-problemphyxmini0659-spectralcolor}
  The observed blue emission line. Its endpoints and wavelength are independent physical data; the governing law below relates them to the level energies
\end{definition}
\begin{proof}
  This is the declared model definition of Mercury Emission Line.
\end{proof}
\begin{definition}[Mercury Excitation Setup]
  \label{def:physics:phyx-mini-0659:phyxminiproblems-problemphyxmini0659-mercuryexcitationsetup}
  \lean{PhyXMiniProblems.ProblemPhyXMini0659.MercuryExcitationSetup}
  \uses{def:physics:phyx-mini-0659:phyxminiproblems-problemphyxmini0659-massquantity, def:physics:phyx-mini-0659:phyxminiproblems-problemphyxmini0659-actionquantity, def:physics:phyx-mini-0659:phyxminiproblems-problemphyxmini0659-lightspeedquantity, def:physics:phyx-mini-0659:phyxminiproblems-problemphyxmini0659-atomicspecies, def:physics:phyx-mini-0659:phyxminiproblems-problemphyxmini0659-projectilespecies, def:physics:phyx-mini-0659:phyxminiproblems-problemphyxmini0659-mercuryenergylevel, def:physics:phyx-mini-0659:phyxminiproblems-problemphyxmini0659-mercuryenergyleveldiagram, def:physics:phyx-mini-0659:phyxminiproblems-problemphyxmini0659-mercuryemissionline}
  Independent quantities in the collision-and-emission experiment. In particular, neither minimumIncidentSpeed nor thresholdKineticEnergy is defined from an answer choice
\end{definition}
\begin{proof}
  This is the declared model definition of Mercury Excitation Setup.
\end{proof}
\begin{definition}[Agrees With Nearest Nanometer Readout]
  \label{def:physics:phyx-mini-0659:phyxminiproblems-problemphyxmini0659-agreeswithnearestnanometerreadout}
  \lean{PhyXMiniProblems.ProblemPhyXMini0659.AgreesWithNearestNanometerReadout}
  \uses{def:physics:phyx-mini-0659:phyxminiproblems-problemphyxmini0659-lengthquantity, def:physics:phyx-mini-0659:phyxminiproblems-problemphyxmini0659-lengthinnanometers}
  Agreement with a wavelength printed to the nearest nanometre
\end{definition}
\begin{proof}
  This is the declared model definition of Agrees With Nearest Nanometer Readout.
\end{proof}
\begin{definition}[Agrees With Two Decimal Electron Volt Readout]
  \label{def:physics:phyx-mini-0659:phyxminiproblems-problemphyxmini0659-agreeswithtwodecimalelectronvoltreadout}
  \lean{PhyXMiniProblems.ProblemPhyXMini0659.AgreesWithTwoDecimalElectronVoltReadout}
  \uses{def:physics:phyx-mini-0659:phyxminiproblems-problemphyxmini0659-energyinelectronvolts}
  Agreement with an energy printed to two decimal places in electron volts
\end{definition}
\begin{proof}
  This is the declared model definition of Agrees With Two Decimal Electron Volt Readout.
\end{proof}
\begin{definition}[Matches Mercury Electron Excitation Scenario]
  \label{def:physics:phyx-mini-0659:phyxminiproblems-problemphyxmini0659-matchesmercuryelectronexcitationscenario}
  \lean{PhyXMiniProblems.ProblemPhyXMini0659.MatchesMercuryElectronExcitationScenario}
  \uses{def:physics:phyx-mini-0659:phyxminiproblems-problemphyxmini0659-mercuryexcitationsetup, def:physics:phyx-mini-0659:phyxminiproblems-problemphyxmini0659-agreeswithnearestnanometerreadout}
  The prose scenario and the independently supplied 492 nm blue-line datum
\end{definition}
\begin{proof}
  This is the declared model definition of Matches Mercury Electron Excitation Scenario.
\end{proof}
\begin{definition}[Matches Supplied Mercury Energy Level Diagram]
  \label{def:physics:phyx-mini-0659:phyxminiproblems-problemphyxmini0659-matchessuppliedmercuryenergyleveldiagram}
  \lean{PhyXMiniProblems.ProblemPhyXMini0659.MatchesSuppliedMercuryEnergyLevelDiagram}
  \uses{def:physics:phyx-mini-0659:phyxminiproblems-problemphyxmini0659-mercuryexcitationsetup, def:physics:phyx-mini-0659:phyxminiproblems-problemphyxmini0659-agreeswithtwodecimalelectronvoltreadout}
  Primary-image evidence from 659.png, including both distinct states printed at 8.84 eV. The approximate physical-energy clauses interpret the two-place labels as measurements and contain no spectral-line endpoint or speed answer
\end{definition}
\begin{proof}
  This is the declared model definition of Matches Supplied Mercury Energy Level Diagram.
\end{proof}
\begin{definition}[Uses Standard Electron And Photon Reference Data]
  \label{def:physics:phyx-mini-0659:phyxminiproblems-problemphyxmini0659-usesstandardelectronandphotonreferencedata}
  \lean{PhyXMiniProblems.ProblemPhyXMini0659.UsesStandardElectronAndPhotonReferenceData}
  \uses{def:physics:phyx-mini-0659:phyxminiproblems-problemphyxmini0659-massinkilograms, def:physics:phyx-mini-0659:phyxminiproblems-problemphyxmini0659-actioninjouleseconds, def:physics:phyx-mini-0659:phyxminiproblems-problemphyxmini0659-lightspeedinmeterspersecond, def:physics:phyx-mini-0659:phyxminiproblems-problemphyxmini0659-mercuryexcitationsetup}
  Standard electron mass, ordinary Planck action, and vacuum light speed
\end{definition}
\begin{proof}
  This is the declared model definition of Uses Standard Electron And Photon Reference Data.
\end{proof}
\begin{definition}[Has Physical Mercury Excitation Parameters]
  \label{def:physics:phyx-mini-0659:phyxminiproblems-problemphyxmini0659-hasphysicalmercuryexcitationparameters}
  \lean{PhyXMiniProblems.ProblemPhyXMini0659.HasPhysicalMercuryExcitationParameters}
  \uses{def:physics:phyx-mini-0659:phyxminiproblems-problemphyxmini0659-lengthinmeters, def:physics:phyx-mini-0659:phyxminiproblems-problemphyxmini0659-massinkilograms, def:physics:phyx-mini-0659:phyxminiproblems-problemphyxmini0659-energyinjoules, def:physics:phyx-mini-0659:phyxminiproblems-problemphyxmini0659-actioninjouleseconds, def:physics:phyx-mini-0659:phyxminiproblems-problemphyxmini0659-speedinmeterspersecond, def:physics:phyx-mini-0659:phyxminiproblems-problemphyxmini0659-lightspeedinmeterspersecond, def:physics:phyx-mini-0659:phyxminiproblems-problemphyxmini0659-mercuryexcitationsetup}
  Positivity and nonnegativity conditions selecting the physical branch
\end{definition}
\begin{proof}
  This is the declared model definition of Has Physical Mercury Excitation Parameters.
\end{proof}
\begin{definition}[target Excitation Energy In Joules]
  \label{def:physics:phyx-mini-0659:phyxminiproblems-problemphyxmini0659-targetexcitationenergyinjoules}
  \lean{PhyXMiniProblems.ProblemPhyXMini0659.targetExcitationEnergyInJoules}
  \uses{def:physics:phyx-mini-0659:phyxminiproblems-problemphyxmini0659-energyinjoules, def:physics:phyx-mini-0659:phyxminiproblems-problemphyxmini0659-mercuryexcitationsetup}
  Coherent-SI energy which the collision must supply to reach the target line's upper level from the atom's initial level
\end{definition}
\begin{proof}
  This is the declared model definition of target Excitation Energy In Joules.
\end{proof}
\begin{definition}[Can Excite Target Line At Speed]
  \label{def:physics:phyx-mini-0659:phyxminiproblems-problemphyxmini0659-canexcitetargetlineatspeed}
  \lean{PhyXMiniProblems.ProblemPhyXMini0659.CanExciteTargetLineAtSpeed}
  \uses{def:physics:phyx-mini-0659:phyxminiproblems-problemphyxmini0659-massinkilograms, def:physics:phyx-mini-0659:phyxminiproblems-problemphyxmini0659-mercuryexcitationsetup, def:physics:phyx-mini-0659:phyxminiproblems-problemphyxmini0659-targetexcitationenergyinjoules}
  At a candidate nonnegative scalar speed, the incident electron can supply the target excitation energy according to nonrelativistic kinetic energy
\end{definition}
\begin{proof}
  This is the declared model definition of Can Excite Target Line At Speed.
\end{proof}
\begin{definition}[Satisfies Mercury Emission And Threshold Laws]
  \label{def:physics:phyx-mini-0659:phyxminiproblems-problemphyxmini0659-satisfiesmercuryemissionandthresholdlaws}
  \lean{PhyXMiniProblems.ProblemPhyXMini0659.SatisfiesMercuryEmissionAndThresholdLaws}
  \uses{def:physics:phyx-mini-0659:phyxminiproblems-problemphyxmini0659-lengthinmeters, def:physics:phyx-mini-0659:phyxminiproblems-problemphyxmini0659-massinkilograms, def:physics:phyx-mini-0659:phyxminiproblems-problemphyxmini0659-energyinjoules, def:physics:phyx-mini-0659:phyxminiproblems-problemphyxmini0659-actioninjouleseconds, def:physics:phyx-mini-0659:phyxminiproblems-problemphyxmini0659-speedinmeterspersecond, def:physics:phyx-mini-0659:phyxminiproblems-problemphyxmini0659-lightspeedinmeterspersecond, def:physics:phyx-mini-0659:phyxminiproblems-problemphyxmini0659-mercuryexcitationsetup, def:physics:phyx-mini-0659:phyxminiproblems-problemphyxmini0659-targetexcitationenergyinjoules, def:physics:phyx-mini-0659:phyxminiproblems-problemphyxmini0659-canexcitetargetlineatspeed}
  General laws used by the problem: the emitted photon's energy gap obeys ΔE λ = h c; the threshold electron has kinetic energy m v² / 2; energy conservation equates that kinetic energy to the excitation from the initial atomic level to the emission line's upper level; and the designated threshold speed is the least nonnegative speed capable of supplying that excitation energy. No clause identifies the 492 nm line's endpoint, assigns a numerical speed, or selects an answer choice
\end{definition}
\begin{proof}
  This is the declared model definition of Satisfies Mercury Emission And Threshold Laws.
\end{proof}
\begin{lemma}[blue492nm Line is level6s8s to level6s6p]
  \label{lem:physics:phyx-mini-0659:phyxminiproblems-problemphyxmini0659-blue492nmline-is-level6s8s-to-level6s6p}
  \lean{PhyXMiniProblems.ProblemPhyXMini0659.blue492nmLine_is_level6s8s_to_level6s6p}
  \uses{def:physics:phyx-mini-0659:phyxminiproblems-problemphyxmini0659-mercuryexcitationsetup, def:physics:phyx-mini-0659:phyxminiproblems-problemphyxmini0659-matchesmercuryelectronexcitationscenario, def:physics:phyx-mini-0659:phyxminiproblems-problemphyxmini0659-matchessuppliedmercuryenergyleveldiagram, def:physics:phyx-mini-0659:phyxminiproblems-problemphyxmini0659-usesstandardelectronandphotonreferencedata, def:physics:phyx-mini-0659:phyxminiproblems-problemphyxmini0659-hasphysicalmercuryexcitationparameters, def:physics:phyx-mini-0659:phyxminiproblems-problemphyxmini0659-satisfiesmercuryemissionandthresholdlaws}
  The only downward pair in the supplied list compatible with the rounded 492 nm wavelength is the 6s8s → 6s6p transition. This is a derived conclusion, not a figure premise
\end{lemma}
\begin{proof}
  The conclusion follows from the governing relations and figure readouts named in the statement of blue492nm Line is level6s8s to level6s6p.
\end{proof}
\begin{lemma}[threshold Energy is ground6s2 to level6s8s]
  \label{lem:physics:phyx-mini-0659:phyxminiproblems-problemphyxmini0659-thresholdenergy-is-ground6s2-to-level6s8s}
  \lean{PhyXMiniProblems.ProblemPhyXMini0659.thresholdEnergy_is_ground6s2_to_level6s8s}
  \uses{def:physics:phyx-mini-0659:phyxminiproblems-problemphyxmini0659-energyinjoules, def:physics:phyx-mini-0659:phyxminiproblems-problemphyxmini0659-mercuryexcitationsetup, def:physics:phyx-mini-0659:phyxminiproblems-problemphyxmini0659-matchesmercuryelectronexcitationscenario, def:physics:phyx-mini-0659:phyxminiproblems-problemphyxmini0659-matchessuppliedmercuryenergyleveldiagram, def:physics:phyx-mini-0659:phyxminiproblems-problemphyxmini0659-usesstandardelectronandphotonreferencedata, def:physics:phyx-mini-0659:phyxminiproblems-problemphyxmini0659-hasphysicalmercuryexcitationparameters, def:physics:phyx-mini-0659:phyxminiproblems-problemphyxmini0659-satisfiesmercuryemissionandthresholdlaws}
  Once the line is identified, threshold energy conservation uses the full ground-to-6s8s excitation gap, not merely the emitted 6s8s → 6s6p gap
\end{lemma}
\begin{proof}
  The conclusion follows from the governing relations and figure readouts named in the statement of threshold Energy is ground6s2 to level6s8s.
\end{proof}
\begin{definition}[Answer Choice]
  \label{def:physics:phyx-mini-0659:phyxminiproblems-problemphyxmini0659-answerchoice}
  \lean{PhyXMiniProblems.ProblemPhyXMini0659.AnswerChoice}
  Labels of the four speed choices supplied with the question
\end{definition}
\begin{proof}
  This is the declared model definition of Answer Choice.
\end{proof}
\begin{definition}[speed Meters Per Second]
  \label{def:physics:phyx-mini-0659:phyxminiproblems-problemphyxmini0659-answerchoice-speedmeterspersecond}
  \lean{PhyXMiniProblems.ProblemPhyXMini0659.AnswerChoice.speedMetersPerSecond}
  \uses{def:physics:phyx-mini-0659:phyxminiproblems-problemphyxmini0659-answerchoice}
  Displayed answer speed in metres per second
\end{definition}
\begin{proof}
  This is the declared model definition of speed Meters Per Second.
\end{proof}
\begin{definition}[recorded Dataset Answer]
  \label{def:physics:phyx-mini-0659:phyxminiproblems-problemphyxmini0659-recordeddatasetanswer}
  \lean{PhyXMiniProblems.ProblemPhyXMini0659.recordedDatasetAnswer}
  \uses{def:physics:phyx-mini-0659:phyxminiproblems-problemphyxmini0659-answerchoice}
  Dataset metadata recording answer C; this is not a theorem premise
\end{definition}
\begin{proof}
  This is the declared model definition of recorded Dataset Answer.
\end{proof}
\begin{definition}[answer Choice Error Meters Per Second]
  \label{def:physics:phyx-mini-0659:phyxminiproblems-problemphyxmini0659-answerchoiceerrormeterspersecond}
  \lean{PhyXMiniProblems.ProblemPhyXMini0659.answerChoiceErrorMetersPerSecond}
  \uses{def:physics:phyx-mini-0659:phyxminiproblems-problemphyxmini0659-speedinmeterspersecond, def:physics:phyx-mini-0659:phyxminiproblems-problemphyxmini0659-mercuryexcitationsetup, def:physics:phyx-mini-0659:phyxminiproblems-problemphyxmini0659-answerchoice}
  Absolute error between the modeled minimum speed and one displayed choice
\end{definition}
\begin{proof}
  This is the declared model definition of answer Choice Error Meters Per Second.
\end{proof}
\begin{definition}[Agrees With Displayed Minimum Speed]
  \label{def:physics:phyx-mini-0659:phyxminiproblems-problemphyxmini0659-agreeswithdisplayedminimumspeed}
  \lean{PhyXMiniProblems.ProblemPhyXMini0659.AgreesWithDisplayedMinimumSpeed}
  \uses{def:physics:phyx-mini-0659:phyxminiproblems-problemphyxmini0659-mercuryexcitationsetup, def:physics:phyx-mini-0659:phyxminiproblems-problemphyxmini0659-answerchoice, def:physics:phyx-mini-0659:phyxminiproblems-problemphyxmini0659-answerchoiceerrormeterspersecond}
  Agreement with a speed whose coefficient is displayed to the nearest hundredth in units of 10\textasciicircum{}6 m/s; half a display unit is 5 × 10\textasciicircum{}3 m/s
\end{definition}
\begin{proof}
  This is the declared model definition of Agrees With Displayed Minimum Speed.
\end{proof}
\begin{definition}[Is Unique Closest Minimum Speed Choice]
  \label{def:physics:phyx-mini-0659:phyxminiproblems-problemphyxmini0659-isuniqueclosestminimumspeedchoice}
  \lean{PhyXMiniProblems.ProblemPhyXMini0659.IsUniqueClosestMinimumSpeedChoice}
  \uses{def:physics:phyx-mini-0659:phyxminiproblems-problemphyxmini0659-mercuryexcitationsetup, def:physics:phyx-mini-0659:phyxminiproblems-problemphyxmini0659-answerchoice, def:physics:phyx-mini-0659:phyxminiproblems-problemphyxmini0659-answerchoiceerrormeterspersecond}
  A choice is strictly closer to the modeled threshold speed than all others
\end{definition}
\begin{proof}
  This is the declared model definition of Is Unique Closest Minimum Speed Choice.
\end{proof}
% --- Archon declaration topology end ---
% --- Archon physics formalization source end ---
